\section{Problem Setup}\label{sec:problem}

The long-only mean-variance portfolio solves
\[
  \min_{w}\; w^\top T_0^{\mathrm{RMT}} w - \rho\,\hat\mu^\top w
  \quad\text{subject to}\quad
  \mathbf{1}^\top w = 1,\quad w \geq 0\,,
\]
where $\rho \geq 0$ is a return-tilt weight, $\hat\mu$ is a vector of expected
returns, and $T_0^{\mathrm{RMT}}$ is the RMT-cleaned covariance estimator
constructed in the previous section (Section~\ref{sec:rmt}).\footnote{The symbol $T_0$ for the cleaned estimator
follows standard shrinkage-target notation; it is unrelated to the sample size $T$.}

\begin{remark}[Why not blend with the sample covariance?]\label{rem:blend}
A general shrinkage estimator
$\hat\Sigma_\alpha = (1-\alpha)\hat\Sigma + \alpha T_0^{\mathrm{RMT}}$ with
$\alpha \in (0,1)$ would re-introduce a fraction of the bulk eigenvalues that
the cleaning replaces. Our reason for fixing $\alpha = 1$ is algebraic, not
statistical: only at $\alpha = 1$ does the system matrix retain the
low-rank-plus-identity form that the Woodbury solve of
Section~\ref{sec:algorithm} requires. For $\alpha \in (0,1)$ the matrix
$(1-\alpha)\hat\Sigma_\mathcal{A} + \alpha T_{0,\mathcal{A}}^{\mathrm{RMT}}$
has full-rank residual structure (no $k \times k$ correction matrix exists)
and the $O(n_a k^2)$ cost advantage is entirely lost. The statistical case for
CRE cleaning itself rests on the RMT literature
(Section~\ref{sec:rmt}); Section~\ref{ssec:oos} reports out-of-sample
evidence that the cleaned estimator is competitive with standard shrinkage on
the data used in this paper.
\end{remark}

The budget constraint $\mathbf{1}^\top w = 1$ is eliminated analytically and the
non-negativity constraint $w \geq 0$ is enforced by a primal-dual active-set
outer loop, each step of which re-solves a balance-constrained subproblem on the
current active set. Section~\ref{sec:algorithm} develops this reduction together
with the Woodbury inner solve that is this paper's contribution.
